%% kapa: exam=1 round=31 year=2020 period=2 form=A subject=accounting subject_ko=회계학 seq=1 number=41
%% source: Q-Net 감정평가사 기출문제 (cst00309) artlSeq=5207219 file=extracted/1008/5207219/A형/A형/2020년도 제31회 감정평가사 1차 시험_2교시_A형.pdf
%% conversion: kordoc 4.24.51 -> tools/gen1.py (auto); review=proofread-sonnet; flags=none
\begin{question}{41}
재무보고를 위한 개념체계 중 재무정보의 질적특성에 관한 설명으로 옳지 않은 것은?
\begin{choices}
  \item 유용한 재무정보의 질적특성은 그 밖의 방법으로 제공되는 재무정보 뿐만 아니라 재무제표에서 제공되는 재무정보에도 적용된다.
  \item 중요성은 기업 특유 관점의 목적적합성을 의미하므로 회계기준위원회는 중요성에 대한 획일적인 계량 임계치를 정하거나 특정한 상황에서 무엇이 중요한 것인지를 미리 결정하여야 한다.
  \item 재무정보의 예측가치와 확인가치는 상호 연관되어 있다. 예측가치를 갖는 정보는 확인가치도 갖는 경우가 많다.
  \item 재무보고의 목적을 달성하기 위해 근본적 질적특성 간 절충(‘trade-off’)이 필요할 수도 있다.
  \item 근본적 질적특성을 충족하면 어느 정도의 비교가능성은 달성될 수 있다.
\end{choices}
\end{question}
\answer{2}
